\documentclass[11pt]{article}
\usepackage[a4paper,margin=2cm]{geometry}
\usepackage{amsmath,amssymb}
\usepackage[hidelinks]{hyperref}
\allowdisplaybreaks
\title{Qpu.Application.Grover}
\author{Tsvetan Rouschev}
\date{@uuidna/qpu v1.1.1, CC-BY-NC-ND-4.0}
\begin{document}
\maketitle
\noindent\url{https://github.com/uuidna/qpu/blob/v1.1.1/lean/Qpu/Application/Grover.lean}

\paragraph{\texttt{groverSize\_growth}}
Unlimited growth: the database for k+1 qubit pairs is exactly four times the k-pair one, for every k — no bound is hidden in a spot value.
\[ \forall (k : \mathbb{N}), \mathrm{groverSize} (k + 1) = 4 \cdot \mathrm{groverSize}\ k \]
\noindent axioms: \texttt{propext}, \texttt{Quot.sound}; proof: \url{https://github.com/uuidna/qpu/blob/v1.1.1/lean/Qpu/Application/Grover.lean\#L36}

\paragraph{\texttt{groverSize\_pos}}
For every k the database is nonempty: there is always something to search.
\[ \forall (k : \mathbb{N}), 0 < \mathrm{groverSize}\ k \]
\noindent axioms: none; proof: \url{https://github.com/uuidna/qpu/blob/v1.1.1/lean/Qpu/Application/Grover.lean\#L41}

\paragraph{\texttt{groverSize\_strict\_growth}}
Growth is strict at every size: adding a qubit pair strictly enlarges the database.
\[ \forall (k : \mathbb{N}), \mathrm{groverSize}\ k < \mathrm{groverSize} (k + 1) \]
\noindent axioms: \texttt{propext}, \texttt{Quot.sound}; proof: \url{https://github.com/uuidna/qpu/blob/v1.1.1/lean/Qpu/Application/Grover.lean\#L44}

\paragraph{\texttt{grover\_all}}
The textbook spot values, decided by the kernel as controls on the growth theorems.
\[ \mathrm{groverIters}\ 2 = 3 \land \mathrm{groverIters}\ 5 = 25 \land \mathrm{groverSize}\ 5 = 1024 \]
\noindent axioms: none; proof: \url{https://github.com/uuidna/qpu/blob/v1.1.1/lean/Qpu/Application/Grover.lean\#L48}
\end{document}
